% !TEX root = agr-paper.tex
% =====================================================================
% NOT A DOCUMENT. Packages and styles for agr-paper.tex, which carries
% the \documentclass and the \begin{document}; building this file
% directly stops with "Missing \begin{document}".
%
% What elsarticle actually loads is fontenc[T1], graphicx, natbib and
% geometry, and nothing else worth relying on. It does NOT load hyperref,
% amsmath, amssymb or url -- \url is undefined until something here
% provides it, which is how the data-availability statement first failed
% to compile. Do not load geometry; elsarticle sets the measure per class
% option and a second geometry undoes it.
% =====================================================================

% Not cosmetic. elsarticle typesets \elsaddress with \footnotesize\itshape
% inside its frontmatter, and under `review' that call reaches font
% selection with an empty size -- "Calculating math sizes for size <>",
% then nine "Font shape OT1/cmr/m/n in size <> not available" warnings
% substituting size <5>. Bisected to \affiliation: remove it and the
% warnings vanish, and all three affiliation forms (full key-value,
% organization-only, plain text) trigger it equally. A scalable family
% clears every one of them, because there is then a shape to find at
% whatever size is asked for. Elsevier re-typesets in their own fonts at
% proof stage, so this affects the submission PDF only. The class already
% sets T1 for text; the warnings are OT1/cmr, which is where maths lives,
% and lmodern is scalable in both.
\usepackage{lmodern}

\usepackage{booktabs}
\usepackage{array}
\usepackage{tikz}
\usepackage{pgfplots}
\pgfplotsset{compat=1.18}
\usetikzlibrary{arrows.meta,positioning,shapes.geometric,calc,fit,
  backgrounds,patterns}

% Semantic figure palette, identical to buetcsepgthesis.sty and to the
% deck: blue for a step that does work, green for supported, vermillion
% for unsupported. The generated figures repeat their own \definecolor
% so they stay self-contained; these are for the hand-drawn ones.
\definecolor{agrNode}{HTML}{0072B2}
\definecolor{agrSup}{HTML}{009E73}
\definecolor{agrUns}{HTML}{D55E00}

% The plot style the generated figures expect. Width and height are set
% per figure by scripts/build_figures.py --target paper, not here: the
% paper target carries its own geometry the way the thesis and slide
% targets do, and repeating it here would silently override nothing
% while misleading the next reader.
\pgfplotsset{
  agrplot/.style={
    tick label style={font=\scriptsize},
    label style={font=\small},
    title style={font=\small\bfseries},
    legend style={font=\scriptsize, draw=none, fill=none},
    grid=major,
    grid style={gray!25, very thin},
    axis line style={gray!60},
    tick align=outside,
    tick style={gray!60},
  }}

% Ragged table cells. Justification spaces badly on a column this narrow,
% and a paper column is narrower than the thesis's.
\newcolumntype{L}[1]{>{\raggedright\arraybackslash}p{#1}}

% Room for the third pass to loosen a line rather than overrun it. The
% names in this literature are long and several are already hyphenated
% ("Think-on-Graph~\cite{...}" set 5.15pt overfull on its own), and a
% 1.5-spaced measure gives TeX fewer break points than the thesis's.
% The class ships no emergencystretch at all. 1em, not more: at 2em the
% third pass starts loosening lines it could have left alone, and the
% run trades its last overfull box for an underfull one.
\emergencystretch=1em

% A preference, not a repair. Breaking at a hyphen this literature
% already contains -- "Reasoning-" / "on-Graphs" -- reads as the start of
% a new word and costs the reader a re-parse, so it is priced high while
% ordinary hyphenation stays cheap. Do not expect it to clear an overfull
% box: swept from 0 to 4000 against the related-work paragraph, the box
% count did not move at any value. Lines that will not fit are fixed by
% rewriting the sentence, which is what happened there.
\hyphenpenalty=400
\exhyphenpenalty=4000

% Short names never break: left to itself TeX splits WebQSP after the
% "We", which a reader has to re-parse. Every entry must be letters only
% -- a digit raises "Not a letter", which is why Neo4j is absent and does
% not need to be here.
\hyphenation{WebQSP CWQ GraphRAG Freebase LangGraph McNemar Cypher AGR}

% ComplexWebQuestions is the exception and must NOT join the list above.
% At twenty letters it is longer than the slack in a 1.5-spaced
% measure, and forbidding every break in it is what set the sentence
% naming both benchmarks 46pt overfull. Hyphens here mark the breaks that
% ARE allowed, so it may split at its own word boundaries and nowhere
% else.
\hyphenation{Complex-Web-Questions}

% Consistent names for the systems, so a rename touches one line rather
% than forty. The thesis writes these out longhand in every table.
\newcommand{\noret}{no-retrieval}
\newcommand{\vecrag}{Vector-RAG}
\newcommand{\graphrag}{GraphRAG}
\newcommand{\tog}{Think-on-Graph}
\newcommand{\agr}{AGR}

% NO LINE NUMBERS, DELIBERATELY. `lineno' was loaded here and activated
% after \end{frontmatter}; both are gone, and the justification
% originally written in this comment did not survive checking. It claimed
% "Elsevier's own guidance asks for numbered lines at submission", which
% is false: Elsevier's generic guide for authors says that at initial
% submission "there are no strict formatting requirements", and does not
% mention line numbering at all (checked August 2026). The KBS guide
% specifically is still unread -- it 403s to an automated fetch by every
% route -- so a journal-level requirement cannot be ruled out; see
% thesis_paper/README.md. Do not re-add on the strength of the old
% comment, which also asserted, without checking, that Editorial Manager
% builds the reviewer PDF itself.
%
% If a submission checklist ever does ask for them, it is two lines:
% \usepackage{lineno} here, and \linenumbers after \end{frontmatter} in
% agr-paper.tex -- in that order and in those two places. Activating in
% the preamble instead numbers the title, author line and abstract
% heading, each on its own inconsistent count before the body starts.
%
% What `review' actually does is \gdef\@blstr{1.5}: 1.5 line spacing, and
% nothing else. It never provided line numbers, whatever the name suggests.

% Loaded last, which is hyperref's requirement: it redefines a great many
% internals and wants everything else already in place. `hidelinks' keeps
% the coloured boxes out of a manuscript that will be read on paper as
% often as on screen, while \Cref, \ref and \cite stay clickable for a
% reviewer who is navigating a PDF. It also supplies \url, which
% elsarticle does not.
\usepackage[hidelinks]{hyperref}

% \url breaks only at punctuation, and the repository URL in the
% data-availability statement has none in the right place: it set 0.68pt
% overfull in the 1.5-spaced measure, which survives however much
% prose surrounds it. xurl lets a URL break anywhere, which is the only
% thing that helps when the unbreakable run is "agentic-graph-reasoning".
% Load it after hyperref -- it patches hyperref's \url.
\usepackage{xurl}

% \Cref throughout, matching the thesis, which gets cleveref from
% buetcsepgthesis.sty. elsarticle does not load it either. cleveref must
% come after hyperref, and last of all, because it redefines the
% referencing commands every other package may have touched.
\usepackage{cleveref}
